%% CV estilo Harvard (formato de la Oficina de Carreras de Harvard) -- Español
%% Compila con pdfLaTeX en Overleaf (Menú > Compilador: pdfLaTeX), sin archivos adicionales.
\documentclass[10pt,letterpaper]{article}

\usepackage[T1]{fontenc}
\usepackage[utf8]{inputenc}
\usepackage[spanish,es-noshorthands]{babel}
\usepackage{mathptmx}                 % Times, the Harvard-style serif
\usepackage[top=0.38in,bottom=0.3in,left=0.6in,right=0.6in]{geometry}
\usepackage{enumitem}
\usepackage{titlesec}
\usepackage[hidelinks]{hyperref}
\usepackage{microtype}

\pagestyle{empty}
\setlength{\parindent}{0pt}
\setlength{\parskip}{0pt}

% Títulos de sección: negrita en versalitas con una línea debajo
\titleformat{\section}{\bfseries\large\scshape}{}{0pt}{}[\vspace{-0.6em}\rule{\linewidth}{0.5pt}]
\titlespacing*{\section}{0pt}{0.5em}{0.3em}

% \entry{Organization}{Location}{Role}{Dates}
\newcommand{\entry}[4]{%
  \textbf{#1}\hfill #2\par\nopagebreak
  \textit{#3}\hfill \textit{#4}\par\nopagebreak\vspace{0.05em}}

% Viñetas compactas: cada una empieza con un verbo de acción
\setlist[itemize]{leftmargin=1.2em,itemsep=0pt,topsep=0.08em,parsep=0pt,partopsep=0pt,label=\textbullet}

\newcommand{\gap}{\vspace{0.22em}}

\begin{document}

% ---------- Header ----------
\begin{center}
  {\LARGE\bfseries Christian Ferrer}\\[0.2em]
  % TODO: agrega tu teléfono si quieres incluirlo -> +56 9 XXXX XXXX $\cdot$
  Temuco, Chile $\cdot$ \href{mailto:christianferrer.dev@gmail.com}{christianferrer.dev@gmail.com} $\cdot$
  \href{https://christianferrer.me}{christianferrer.me}\\
  \href{https://www.linkedin.com/in/christian-ferrer-640b7b216/}{linkedin.com/in/christian-ferrer-640b7b216} $\cdot$
  \href{https://github.com/Ozauri0}{github.com/Ozauri0}
\end{center}

% ---------- Educación ----------
\section{Educación}
% TODO: confirma la institución, el nombre del programa y la fecha estimada de egreso.
\entry{Universidad Católica de Temuco}{Temuco, Chile}{Ingeniería Civil en Informática}{En curso}

% ---------- Software experience ----------
\section{Experiencia en Desarrollo de Software}
\entry{Temblo}{\href{https://temblo.cl}{temblo.cl}}{Desarrollador Full Stack (proyecto personal)}{Actualidad}
\begin{itemize}
  \item Construyo y opero una plataforma de monitoreo sísmico 3D en tiempo real para Chile y el mundo, que recibe telemetría global y detecta sismos automáticamente en segundos.
  \item Desarrollo el stack completo: interfaz en Next.js, TypeScript y Tailwind CSS (vista 3D, control de cámara, estaciones, eventos archivados, alarma sonora) y una API en Node.js/Express con MongoDB y JWT; desplegada con Docker y Nginx.
\end{itemize}
\gap
\entry{GruasCares}{\href{https://gruascares.cl}{gruascares.cl}}{Desarrollador Full Stack}{Actualidad}
\begin{itemize}
  \item Desarrollo una plataforma full stack para Grúas Cares, empresa de rescate de vehículos y serviteca: landing page, solicitud de grúas, agenda de la serviteca y gestión de personal y trabajos, con paneles por rol (cliente, mecánico, admin).
  \item Implemento autenticación por roles (JWT en cookies httpOnly, límite de peticiones) con Next.js, TypeScript, Express sobre Bun, PostgreSQL y Drizzle ORM; con pruebas y contenedores Docker.
\end{itemize}
\gap
\entry{Universidad Católica de Temuco}{Remoto}{Ayudante de Gestión de Datos y Monitoreo (PACE)}{sept.\ 2025 -- actualidad}
\begin{itemize}
  \item Desarrollé, por iniciativa propia, una aplicación de escritorio con Electron y TypeScript que automatiza la descarga masiva de reportes de cursos de Moodle mediante web scraping, eliminando un proceso manual.
  \item Construí un módulo que filtra los registros de docentes y ayudantes y fusiona los archivos de cada curso en un Excel con una hoja por curso.
\end{itemize}
\gap
\entry{Universidad Católica de Temuco (DGE)}{Remoto}{Especialista en Migración y Procesamiento de Datos Académicos (Python/XML)}{nov.\ 2025 -- dic.\ 2025}
\begin{itemize}
  \item Automaticé la publicación digital de las actas del congreso CLABES en Open Journal Systems (OJS) con dos scripts en Python: uno separa y renombra las actas en artículos y el otro genera metadatos XML para la importación masiva, sin intervención manual.
\end{itemize}
\gap
\entry{MiPortalVentas}{}{Desarrollador de Software}{2025}
\begin{itemize}
  \item Desarrollé un sistema ERP para MiPortalVentas.
\end{itemize}

% ---------- Liderazgo ----------
\section{Experiencia de Liderazgo}
\entry{Universidad Católica de Temuco}{Temuco, Chile}{Jefe Coordinador de Procesos Académicos, Programa Bachiller}{jul.\ 2026 -- sept.\ 2026}
\begin{itemize}
  \item Dirigí la operación del área durante el cambio de semestre, a cargo de un equipo de 3 asistentes y reportando directamente a la Jefatura de Carrera.
  \item Revisé y optimicé los horarios de aproximadamente 300 estudiantes, equilibrando oferta, cupos por sección y salas.
  \item Resolví con rapidez contingencias de planificación docente, reestructurando horarios ante cambios imprevistos.
  \item Antes Asistente Coordinador (jul.\ -- ago.\ 2025; ene.\ -- abr.\ 2026): apoyé la matrícula y el control de cupos.
\end{itemize}

% ---------- Additional experience ----------
\section{Otras Experiencias}
\entry{Gobierno Regional de La Araucanía}{Temuco, Chile}{Práctica Inicial, Tecnologías de la Información}{dic.\ 2025 -- ene.\ 2026}
\begin{itemize}
  \item Configuré switches Cisco y telefonía PBX, gestioné usuarios con Active Directory, actualicé hardware e instalé Windows; brindé soporte técnico directo al Consejo Regional y al Gobernador.
\end{itemize}
\gap
\entry{CINAP}{Híbrido}{Tutor E-learning}{dic.\ 2023 -- actualidad}

% ---------- Skills ----------
\section{Habilidades}
\begin{itemize}[label={}, leftmargin=0pt]
  \item \textbf{Lenguajes:} Python, TypeScript, JavaScript
  \item \textbf{Web y backend:} Next.js, React, Angular, Node.js/Express, Django, Flask, FastAPI, REST APIs
  \item \textbf{Datos y DevOps:} MongoDB, MySQL, SQLite, Docker, Linux, Nginx, Git
  \item \textbf{Otros:} Electron, Ionic/Capacitor, web scraping, XML, Cisco IOS, Active Directory
  % TODO: agrega idiomas y nivel, p. ej. \textbf{Idiomas:} Español (nativo), Inglés (...)
\end{itemize}

\end{document}
